\documentclass[11pt]{article}
\usepackage[T1]{fontenc}
\usepackage{lmodern,amsmath,amssymb,amsthm,mathtools,booktabs,microtype}
\usepackage[margin=1in]{geometry}
\usepackage[colorlinks=true,linkcolor=blue,citecolor=blue,urlcolor=blue]{hyperref}
\usepackage{enumitem}
\setlist{itemsep=3pt,topsep=5pt}
\newtheorem{theorem}{Theorem}[section]
\newtheorem{lemma}[theorem]{Lemma}
\newtheorem{proposition}[theorem]{Proposition}
\newtheorem{corollary}[theorem]{Corollary}
\theoremstyle{remark}\newtheorem{remark}[theorem]{Remark}
\newcommand{\R}{\mathbb R}\newcommand{\C}{\mathbb C}
\newcommand{\dd}{\,\mathrm d}\newcommand{\supp}{\operatorname{supp}}
\newcommand{\Var}{\operatorname{Var}}\newcommand{\Res}{\operatorname{Res}}
\newcommand{\Rea}{\operatorname{Re}}\newcommand{\Ima}{\operatorname{Im}}
\newcommand{\norm}[1]{\left\lVert#1\right\rVert}
\newcommand{\abs}[1]{\left\lvert#1\right\rvert}
\title{A sharp cancellation inequality for damped modes\\
and a normal form for shortest positive filters}
\author{Research note prepared for the ProveIt program}
\date{27 November 2026}
\begin{document}\maketitle

\begin{abstract}
A positive filter is a compactly supported probability measure, and cancelling
a damped oscillatory mode means prescribing a zero of its bilateral Laplace
transform. We prove that cancelling $N$ upper-half-plane targets, counted with
multiplicity, requires support diameter at least $N^2\pi/\sum_j\Ima s_j$.
The damping rates may be unrelated. Equality is characterized exactly by a
reflection condition and a polynomial with nonnegative coefficients.

The main result describes the entire local geometry at every equality
configuration having strictly positive coefficients. After fixing the complex
mean of the targets, the coefficients of their centered polynomial provide
coordinates that remain valid at collisions. The shortest width is real
analytic, has a positive-definite Hessian transverse to the reflection locus,
and admits the exact local form $N\pi/\kappa+\norm{\eta}^2$. We give its Hessian
as a weighted complete-graph energy with explicit Hermite interpolation.
This answers the sign and degeneracy part of the repository's open question
about second variations at general mirror configurations. An exact cubic
example shows that mixed-derivative signs can change even though positive
definiteness persists. The theorems are unconditional Fourier--Laplace
analysis; no assertion about a new major number-theoretic conjecture is made.
\end{abstract}

\section{Scope, provenance, and the main result}

The starting point is the section ``Positive filters that cancel zeros'' of
the ProveIt article, pinned to commit
\texttt{cfab41369737c66ae520d179419509eb719367cd} \cite{ProveIt}.
That section proves an exact reflection theorem, an analytic continuation
theorem when target multiplicities are held fixed, a mean-frequency inequality
for targets with a common real part, and a detailed Hessian theorem at a
single coalesced frequency. Its open-question list asks for the Hessian at a
general mirror configuration, especially its sign for perturbations fixing
the mean imaginary part. We resolve this question in coefficient coordinates,
allowing arbitrary damping and allowing multiple targets to split.

The contribution is the general normal form and its explicit positive
quadratic form. The strip identity and the reflection certificate used below
are mechanisms already present in the repository; they are proved here to
make the argument self-contained. The global inequality for unrelated real
parts is a short additional consequence. The local normal form does not
solve the global support-transition problem at a configuration where some
optimal masses are zero.

For a compactly supported probability measure $\mu$ on $\R$, put
\[
 F_\mu(s)=\int_\R e^{st}\dd\mu(t).
\]
A multiset $S=(s_1,\ldots,s_N)$ in $\Ima s>0$ specifies zeros of $F_\mu$,
with repeated entries specifying the corresponding multiplicity. Write
$D(S)$ for the least possible diameter of $\supp\mu$. This quantity is
unchanged by translating the measure. It is also unchanged by subtracting
a common real number from all targets: exponential tilting of the measure
gives the required bijection of feasible measures.

Here is the main result. Its notation is developed in the proof.

\begin{theorem}[Normal form at every positive mirror configuration]
\label{thm:main}
Fix $N\ge2$ and $\kappa>0$, and put $h=\pi/\kappa$. Let
\[
 P_0(y)=y^N+\sum_{r=2}^N c_r^0y^{N-r},\qquad c_r^0\in\R,
\]
have all its roots $y_j$, counted with multiplicity, in
$\abs{\Ima y}<\kappa$. Suppose that
\[
 A(z)=\prod_{j=1}^N(z+e^{hy_j})=\sum_{k=0}^N a_kz^k
 \quad\hbox{has }a_k>0\text{ for every }k.
\]
For complex coefficient vectors $c=(c_2,\ldots,c_N)$ close to $c^0$, let
$S_c$ be the multiset of targets $i\kappa+y$, where $y$ runs over the roots
of $P_c(y)=y^N+\sum_{r=2}^N c_ry^{N-r}$. Then:
\begin{enumerate}[label=(\roman*)]
\item The shortest filter is unique up to translation. With left endpoint
zero, it has exactly $N+1$ strictly positive atoms; its nodes, masses, and
width $D(c)$ are real analytic in $(\Rea c,\Ima c)$, including at repeated
roots of $P_c$.
\item The equality locus is exactly $\Ima c=0$. There are constants
$0<C_1<C_2<\infty$ and a neighborhood of $c^0$ on which
\[
 C_1\norm{\Ima c}^2\le D(c)-Nh\le C_2\norm{\Ima c}^2.
\]
The Hessian has kernel exactly the tangent space to the real coefficient
locus. Its restriction to the imaginary coefficient directions is
positive definite.
\item There is a real analytic change of coordinates
$(\Rea c,\Ima c)\mapsto(\Rea c,\eta)$, preserving that locus, for which
\[
 \boxed{\quad D(c)=Nh+\norm{\eta}^2.\quad}
\]
\item For a real polynomial $V(y)$ of degree at most $N-2$, let
$P_\epsilon=P_0+i\epsilon V$ and write the extremal nodes as
$t_k(\epsilon)=kh+\epsilon\sigma_k+O(\epsilon^2)$. Then
$\sigma_0=\sigma_N=0$ and
\begin{equation}\label{eq:energy}
 D(P_\epsilon)=Nh+\epsilon^2 Q_{P_0}(V)+O(\epsilon^3),\qquad
 Q_{P_0}(V)=\frac1h\sum_{0\le j<k\le N}
 a_ja_k(\sigma_j-\sigma_k)^2
 \int_0^\infty\frac{z^{j+k-1}}{A(z)^2}\dd z.
\end{equation}
The linear map $V\mapsto(\sigma_1,\ldots,\sigma_{N-1})$ is an isomorphism.
An explicit interpolation formula for it is given in
Proposition~\ref{prop:interpolation}.
\end{enumerate}
\end{theorem}

The coefficients, rather than individually labeled roots, are essential.
Near a repeated root a root displacement can change the coefficient vector
only at high order. Thus the theorem is consistent with the sixth-, tenth-,
and higher-order width effects already proved in the article. There is no
contradiction between nondegeneracy in coefficient coordinates and high-order
flatness along special root-coordinate paths.

\section{A sharp global inequality for arbitrary damped targets}

\subsection{Existence and the strip identity}

For a single target $\beta+i\gamma$ with $\gamma>0$, a positive two-point
measure on $0,\pi/\gamma$, with masses in ratio
$e^{\beta\pi/\gamma}:1$, has the required simple zero. Convolution gives
feasible measures for every finite multiset and every multiplicity. After
translation, a minimizing sequence lies in one fixed compact interval;
weak compactness, and continuity of every jet test $t^ke^{st}$ on that
interval, give attainment of $D(S)$.

\begin{lemma}[Finite exponential-sum strip identity]\label{lem:strip}
Let $0=t_0<t_1<\cdots<t_M=L$ and let the positive numbers $w_j$ have sum
one. Put $G(s)=\sum_jw_je^{t_js}$. Then, for $T>0$,
\begin{equation}\label{eq:strip}
 J(T):=\int_\R\log\frac{G(x)}{\abs{G(x+iT)}}\dd x
 =\frac{LT^2}{2}-2\pi\sum_{0<\Ima\rho<T}(T-\Ima\rho),
 \qquad J(T)\ge0,
\end{equation}
where zeros are counted with multiplicity. Zeros on the top boundary are
understood by continuity and contribute zero to the sum.
\end{lemma}

\begin{proof}
On the right and left ends of the real axis, the last and first terms,
respectively, dominate uniformly in the imaginary part. Thus all zeros have
bounded real part. The logarithm in the integral decays exponentially at
both real ends; its singularities at zeros are integrable. The distributional
identity
$\Delta\log\abs{G}=2\pi\sum_\rho\delta_\rho$ follows by factoring each
zero locally. Integrating in $x$, the $x$-derivative term has endpoint
contribution $L$, because the limiting logarithmic slopes are $L$ and $0$.
Consequently
$J''(T)=L-2\pi\sum_\rho\delta_{\Ima\rho}(T)$ for $T>0$.
The initial values are $J(0)=J'(0)=0$: $G$ is positive on the real line
and its coefficients are real. Two integrations give \eqref{eq:strip}.
Finally $\abs{G(x+iT)}\le G(x)$ is the triangle inequality.
\end{proof}

\begin{theorem}[Mean imaginary-part inequality and equality cases]
\label{thm:global}
For arbitrary upper-half-plane targets $S=(s_1,\ldots,s_N)$, counted with
multiplicity, put $\kappa=N^{-1}\sum_j\Ima s_j$ and $h=\pi/\kappa$. Then
\begin{equation}\label{eq:global}
 \boxed{\qquad D(S)\ge \frac{N\pi}{\kappa}
       =\frac{N^2\pi}{\sum_j\Ima s_j}.\qquad}
\end{equation}
Equality holds if and only if:
\begin{enumerate}[label=(\alph*)]
\item $0<\Ima s_j<2\kappa$ for every $j$;
\item the target multiset is invariant under
$s\mapsto\overline{s}+2i\kappa$;
\item $A_S(z)=\prod_j(z-e^{hs_j})$ has nonnegative coefficients.
\end{enumerate}
At equality the unique minimizing measure up to translation has masses
given by the coefficients of $A_S(z)/A_S(1)$ at $0,h,\ldots,Nh$.
\end{theorem}

\begin{proof}
The real and imaginary parts of the jet constraints form a finite-dimensional
continuous vector of functions on any fixed support interval. Its integral
can be realized by a finite positive measure on that interval, by the
finite-dimensional convex hull theorem. Thus an optimal measure can be
replaced by a finite one without increasing its width. Translate its support
to $[0,L]$ and discard any zero masses. Apply Lemma~\ref{lem:strip} with
$T=2\kappa$. Its zero sum is at least
\[
 \sum_{j=1}^N(T-\Ima s_j)_+
 \ge NT-\sum_j\Ima s_j=N\kappa.
\]
Since $J(T)\ge0$, this yields $2\kappa^2L\ge2\pi N\kappa$.

At equality $J(T)=0$, and the continuous nonnegative integrand must vanish
away from its possible singularities; a top-boundary zero would make the
integral strictly positive. Equality in the triangle inequality at $x=0$
shows that every $e^{iTt_j}$ is $1$. Hence the nodes belong to $h\mathbb Z$
and $G(s)=B(e^{hs})$, where $B$ has degree $N$, nonnegative coefficients,
and $B(1)=1$. Equality in the preceding sum requires $\Ima s_j\le2\kappa$;
the equality case at the top boundary is excluded because $B(e^{hs})>0$
there for real $s-2i\kappa$. Exponentiation is injective on the open strip
$0<\Ima s<2\kappa$, so the prescribed roots determine $B=A_S/A_S(1)$,
with their multiplicities. Since $B$ is real, its roots are closed under
conjugation, which gives (b). This proves necessity.

Conversely (a)--(c) produce the stated positive measure, whose transform
has exactly the required zeros and whose endpoint masses are positive.
Its diameter is $Nh$, so the lower bound is attained. Uniqueness for
arbitrary, possibly nonfinite, minimizing measures follows from the
certificate in Lemma~\ref{lem:certificate}: every such measure is supported
on its lattice zero set, after which polynomial divisibility fixes all
masses.
\end{proof}

The inequality does not depend on the real parts of the targets. Its equality
conditions do. It therefore distinguishes a lower bound which is insensitive
to damping from the exact compatibility conditions required to attain it.
The integrated angular inequality of Eremenko and Fryntov
\cite{EremenkoFryntov} is a related statement for polynomials with
nonnegative coefficients; the variable here is support duration rather than
polynomial degree.

\subsection{The reflection certificate}

\begin{lemma}[A certificate with nondegenerate zeros]\label{lem:certificate}
Under conditions (a)--(b) of Theorem~\ref{thm:global}, there is a real
exponential polynomial $g$ in the span of the prescribed real and imaginary
jet functions such that $g\ge0$ on $[0,Nh]$, with exactly the zeros
$0,h,\ldots,Nh$. Its endpoint zeros are simple and its interior zeros are
double with strictly positive second derivatives.
\end{lemma}

\begin{proof}
Put $q_j=-e^{hs_j}$ and $Q(z)=\prod_j(z-q_j)$. The branch
$\log q_j=hs_j-i\pi$ is defined by $-\pi<\Ima\log q_j<\pi$.
Reflection invariance makes the multiset of $q_j$ conjugation invariant;
real $q_j$ are positive. Thus
$Q(-t)=(-1)^N\abs{Q(-t)}$ for $t\ge0$.
With repeated poles included, define
\[
 E(x)=\sum_{q}\Res_{z=q}\frac{z^{x-1}}{Q(z)},\qquad
 g(hx)=(-1)^{N+1}\sin(\pi x)E(x).
\]
The residue formula expresses $E$ in the functions
$x^ke^{x\log q_j}$. Multiplication by $\sin(\pi x)$ gives exactly the
target and conjugate-target jet span. A keyhole contour around the negative
real axis gives, for $0<x<N$,
\begin{equation}\label{eq:certificate}
 g(hx)=\frac{\sin^2(\pi x)}{\pi}
       \int_0^\infty\frac{t^{x-1}}{\abs{Q(-t)}}\dd t.
\end{equation}
The outer and inner contour integrals tend to zero because their bounds
are $O(R^{x-N})$ and $O(\epsilon^x)$, respectively. The integral in
\eqref{eq:certificate} is strictly positive. At $x\downarrow0$ it is
$1/(\abs{Q(0)}x)+O(1)$, and at $x\uparrow N$ it is
$1/(N-x)+O(1)$. These formulas prove the asserted simple endpoint zeros,
double interior zeros, and all strict signs.
\end{proof}

\section{Coordinates that remain regular at multiple targets}

Fix the data of Theorem~\ref{thm:main}. The constraint $c_1=0$ fixes the
complex mean of the targets at $i\kappa$. No generality is lost locally:
first remove the common real mean by exponential tilting and then fix the
positive imaginary mean. In these coordinates a multiset is a mirror
configuration exactly when its centered polynomial has real coefficients.

Choose a positively oriented contour $C$ surrounding the roots of $P_0$,
and no nearby root lies on $C$. For an entire function $f$, define
\begin{equation}\label{eq:remainder}
 (\mathcal R_cf)(y)=\frac1{2\pi i}\int_C
 \frac{f(z)}{P_c(z)}\frac{P_c(z)-P_c(y)}{z-y}\dd z.
\end{equation}
This is the polynomial of degree below $N$ which agrees with $f$ and its
required derivatives at the roots of $P_c$. It depends holomorphically on
$c$. Indeed Cauchy's formula rewrites it inside $C$ as $f(y)$ minus
$P_c(y)$ times a holomorphic function. Consequently
\[
 \mathcal R_c\bigl[F_\mu(i\kappa+\,\cdot\,)\bigr]=0
\]
is equivalent to all zero and jet conditions, even at collisions.

\begin{proposition}[Analytic optimizer in coefficient space]
\label{prop:analytic}
For $c$ close to $c^0$, there are real analytic positive weights $w_k(c)$
and increasing nodes $0=t_0(c)<\cdots<t_N(c)$ satisfying the constraints,
with $t_k(c^0)=kh$ and $w_k(c^0)=a_k/A(1)$. This measure is the unique
global optimizer after anchoring its left endpoint at zero.
\end{proposition}

\begin{proof}
There are $2N+1$ real unknowns: $N+1$ weights and $N$ nonzero nodes. The
equations are total mass one and the real and imaginary parts of the $N$
coefficients of \eqref{eq:remainder}. To check the Jacobian at the base,
write a homogeneous variation of the transform as
\[
 \dot F(s)=W(e^{hs})+sB(e^{hs}),\quad
 W(z)=\sum_k\dot w_kz^k,\quad B(z)=\sum_kw_k\dot t_kz^k.
\]
Both polynomials are real, $W(1)=0$, and $B(0)=0$. The homogeneous
constraint says
\[
 P_0(y)\mid W(-e^{hy})+(y+i\kappa)B(-e^{hy}).
\]
Here divisibility means vanishing with the required multiplicities, or
equivalently zero remainder. Since $P_0$ is real, conjugating the
coefficients of the analytic function in $y$ gives the same statement with
$i\kappa$ replaced by $-i\kappa$. Subtraction shows that
$P_0(y)\mid B(-e^{hy})$. The exponential map is injective on the root strip
and has nonzero derivative, so $A\mid B$. Since $\deg B\le N$ and
$B(0)=0\ne A(0)$, $B=0$. Then $A\mid W$, and $W(1)=0\ne A(1)$ gives
$W=0$. All $w_k>0$, so all node and weight variations vanish. The real
Jacobian is invertible, and the real analytic implicit function theorem
produces the asserted branch. In particular this argument does not require
distinct roots.

For completeness, global optimality follows from a persistent dual
certificate, not merely from solving stationary equations. Let $Y_c(t)$
be the $2N$ real functions obtained by taking real and imaginary
coefficients of $\mathcal R_c[e^{(i\kappa+y)t}]$. The vectors
$Y_c(t_k)$ for $0\le k\le N$ and $Y_c'(t_k)$ for $1\le k<N$ span a
hyperplane. They have the positive weight dependence, and adding
$Y_c'(t_N)$ gives full rank by the Jacobian calculation. Choose the normal
analytically with a fixed nonzero normalization on $Y_c'(t_N)$.
Its scalar product with $Y_c(t)$ is a certificate vanishing at all nodes
and with vanishing derivative at interior nodes. At the base it is a
positive multiple of Lemma~\ref{lem:certificate}'s certificate. Strict
endpoint slopes, positive interior curvatures, and positivity away from
small neighborhoods of these finitely many zeros persist. Thus the
certificate is nonnegative on the entire new interval and has exactly
the new nodes as zeros.

Every feasible measure of diameter no greater than the new width can be
translated into this interval and must live on those nodes. The mass row
together with the moment columns is independent, again by the full
Jacobian, so the weights are uniquely determined. Both endpoint weights
are positive, excluding a smaller diameter.
\end{proof}

\section{The exact deficit and the general Hessian}

\subsection{No unaccounted zeros in the strip}

\begin{lemma}[Exact deficit identity near the mirror locus]
\label{lem:deficit}
Write $G_c(s)=\sum_kw_k(c)e^{t_k(c)s}$ and
\[
 p_k(x,c)=\frac{w_k(c)e^{xt_k(c)}}{G_c(x)}.
\]
After shrinking the coefficient neighborhood, $G_c$ has exactly $N$ zeros
in $0<\Ima s<2\kappa$, namely the prescribed targets. Moreover
\begin{equation}\label{eq:deficit}
 D(c)-Nh=\frac1{2\kappa^2}\int_\R
 -\log\abs{\sum_{k=0}^N p_k(x,c)e^{2i\kappa t_k(c)}}\dd x.
\end{equation}
\end{lemma}

\begin{proof}
The endpoint weights have a uniform positive lower bound and all successive
node gaps have a uniform positive lower bound on a sufficiently small
closed coefficient neighborhood. On $\Rea s>R$, division by the last term
bounds the sum of the other absolute values by $Ce^{-gR}$; on
$\Rea s<-R$, use the first term. Choosing $R$ large therefore excludes
zeros outside one fixed vertical rectangle, uniformly in $c$ and in
$0\le\Ima s\le2\kappa$.

At the base $G_0(s)=A(e^{hs})/A(1)$. It has precisely $N$ zeros in that
strip, since each polynomial root has exactly one logarithm there.
Neither horizontal boundary contains a zero: $G_0(x)$ and
$G_0(x+2i\kappa)$ are the same positive number. Uniform convergence on
the compact rectangle and the argument principle preserve the number of
zeros. The prescribed target polynomial supplies $N$ of them, counted
with multiplicity, and all remain inside the strip. Thus there are no
extra zeros. Their imaginary parts sum to $N\kappa$ because $c_1=0$.
Lemma~\ref{lem:strip} at $T=2\kappa$ now gives
$J(2\kappa)=2\kappa^2(D(c)-Nh)$. Dividing its integrand by $G_c(x)$
gives \eqref{eq:deficit}.
\end{proof}

The identity is useful because it expresses the width loss as an integral
of failures of equality in the triangle inequality. It has no unspecified
zero contribution. Endpoint dominance is what makes that assertion uniform
when targets split or collide.

\subsection{The quadratic form}

\begin{proposition}[Variance and graph-energy formulas]
\label{prop:variance}
For any real analytic coefficient path through the base with node
velocities $\sigma_k$, the coefficient of $\epsilon^2$ in $D-Nh$ is
\begin{equation}\label{eq:variance}
 Q=\int_\R \Var_{p(x,c^0)}(\sigma_k)\dd x.
\end{equation}
Equivalently it is the right side of \eqref{eq:energy}.
\end{proposition}

\begin{proof}
At the base $e^{2i\kappa kh}=1$. For any probability vector $p$ and real
$\sigma_k$,
\[
 -\log\abs{\sum_kp_ke^{2i\kappa\epsilon\sigma_k}}
 =2\kappa^2\epsilon^2\Var_p(\sigma_k)+O(\epsilon^3).
\]
Changes of $p$ do not contribute to the quadratic coefficient: its total
mass remains one, the first-order modulus loss is zero, and replacing $p$
by its base value in the displayed second-order term changes only higher
orders. Differentiation under the integral in \eqref{eq:deficit} is valid.
Indeed, after factoring out the endpoint phase, endpoint dominance bounds
the first two parameter derivatives of the integrand by
$C(1+\abs x)^2e^{-g\abs x}$ on the tails. On the remaining compact
interval the phase sum is bounded away from zero near the base.

Now set $z=e^{hx}$. Then $\dd x=\dd z/(hz)$ and
$p_k=a_kz^k/A(z)$. The identity
\[
 \Var_p(\sigma)=\sum_{j<k}p_jp_k(\sigma_j-\sigma_k)^2
\]
gives \eqref{eq:energy}. Every integral is finite: near zero the exponent
is at least zero, and near infinity the largest exponent after division
by $A^2$ is $-2$.
\end{proof}

\begin{proposition}[Explicit velocity map and its invertibility]
\label{prop:interpolation}
For a real $V$ of degree at most $N-2$, let
\[
 H(y)=\frac{A(-e^{hy})}{A(1)P_0(y)}.
\]
The apparent singularities of $H$ at roots of $P_0$ are removable and its
values there are nonzero. The real polynomial
\[
 B_V(z)=\sum_{k=0}^N\frac{a_k}{A(1)}\sigma_k z^k
\]
is characterized by
\begin{equation}\label{eq:hermite}
 \deg B_V\le N-1,\quad B_V(0)=0,\qquad
 \kappa B_V(-e^{hy})\equiv V(y)H(y)\pmod{P_0(y)}.
\end{equation}
It depends linearly and injectively on $V$.
If the $y_j$ are distinct, an explicit formula is
\begin{equation}\label{eq:explicitR}
 R_V(z)=\sum_{j=1}^N\frac{z_jV(y_j)}{P_0'(y_j)}
                     \frac{A(z)}{z-z_j},\quad z_j=-e^{hy_j},\qquad
 \sigma_k=\frac h\kappa\frac{[z^k]R_V(z)}{a_k}.
\end{equation}
For repeated roots, the unique confluent interpolation in
\eqref{eq:hermite}, or the continuous limit of \eqref{eq:explicitR},
defines the same map.
\end{proposition}

\begin{proof}
Along $P_\epsilon=P_0+i\epsilon V$, differentiate the divisibility
$P_\epsilon\mid G_\epsilon(i\kappa+\,\cdot\,)$. At the base this gives
\[
 W(-e^{hy})+(y+i\kappa)B_V(-e^{hy})
       \equiv iV(y)H(y)\pmod{P_0(y)}.
\]
Both $P_0$ and $H$ are real in the sense of conjugating all coefficients.
Taking imaginary parts in the real quotient algebra gives
\eqref{eq:hermite}'s congruence. Also $B_V(0)=0$ by anchoring, and the
global bound has a minimum at the base along this two-sided path, so
$\dot D=0$. This is exactly the vanishing of the leading coefficient
of $B_V$, since $w_N>0$.

If $B_V=0$, then $P_0\mid VH$. The local nonvanishing of $H$ implies
$P_0\mid V$, including multiplicities. Since $\deg V<N$, $V=0$.
Thus the map is injective. Its source and target both have real dimension
$N-1$, so it is an isomorphism onto the internal node velocities.

For distinct roots,
$H(y_j)=hz_jA'(z_j)/(A(1)P_0'(y_j))$. Lagrange interpolation now gives
$B_V=hR_V/(\kappa A(1))$. The constant coefficient vanishes because
\[
 R_V(0)=-A(0)\sum_j\frac{V(y_j)}{P_0'(y_j)}=0;
\]
the last sum is zero for a polynomial of degree at most $N-2$, by the
leading-coefficient formula for interpolation. Finally the maps involved
are analytic in the coefficient data. Simple real coefficient polynomials
are dense among real coefficient polynomials, and positivity and the strip
condition persist in a sufficiently small neighborhood. The formula
therefore extends continuously to repeated roots and agrees with the
unique congruence solution there.
\end{proof}

All edge weights in \eqref{eq:energy} are strictly positive. Its kernel
on vectors $\sigma$ is the constant vectors; anchoring gives
$\sigma_0=0$, so it has zero kernel on internal velocities with fixed
endpoints. Proposition~\ref{prop:interpolation} therefore proves
$Q_{P_0}(V)>0$ for every $V\ne0$. This establishes the asserted transverse
positive definiteness without requiring any spacing assumption on target
roots.

\section{Exact local normal form and quantitative stability}

Write $c=E+iO$ with real coefficient vectors. For all nearby $E$, the
reflection conditions persist and all coefficients of the corresponding
exponential-root polynomial remain positive. Hence $D(E,0)=Nh$.
The global inequality gives $D(E,O)\ge Nh$, so the first derivative in
$O$ vanishes at $O=0$. Real analyticity and Taylor's theorem in $O$ give
\begin{equation}\label{eq:factor}
 D(E,O)-Nh=O^{\mathsf T}M(E,O)O,
\end{equation}
where $M$ is a real analytic symmetric matrix; for example one may take
\[
 M(E,O)=\int_0^1(1-t)\,\partial_{OO}^2D(E,tO)\dd t.
\]
The preceding section proves that $M(c^0,0)$ is positive definite. Thus
all its eigenvalues stay between two fixed positive constants on a
sufficiently small neighborhood. This proves both stability inequalities
and the exact kernel of the full Hessian. The mixed $E,O$ block vanishes
on $O=0$ because $\partial_OD(E,0)=0$, and the $E,E$ block vanishes
because $D(E,0)$ is constant.

A positive definite matrix near $M(c^0,0)$ has a real analytic positive
definite square root. One justification is the analytic implicit function
theorem applied to $X^2=M$: on symmetric matrices the derivative
$Z\mapsto XZ+ZX$ is invertible when $X$ is positive definite, since in
an eigenbasis it multiplies the $(j,k)$ entry by the positive number
$\lambda_j+\lambda_k$. Set
$\eta=M(E,O)^{1/2}O$. The derivative in $O$ at the base is invertible,
so $(E,O)\mapsto(E,\eta)$ is a local real analytic coordinate change.
Equation~\eqref{eq:factor} becomes the exact claimed normal form.
This completes the proof of Theorem~\ref{thm:main}.

\begin{corollary}[Root-coordinate interpretation at simple targets]
At a positive mirror configuration with distinct targets and fixed complex
mean, the second variation vanishes precisely for infinitesimal motions
which preserve the reflection pairing. Every transverse infinitesimal
motion has strictly positive second variation.
\end{corollary}

\begin{proof}
For distinct roots the change from labeled roots to polynomial coefficients
has invertible complex Jacobian after fixing their sum. The polynomial
variation has real coefficients exactly when the root velocities obey
$\dot y_{j'}=\overline{\dot y_j}$ for conjugate partners. Apply the
Hessian kernel statement.
\end{proof}

\begin{remark}[Collisions and unusually high powers]
At a collision the last corollary's root-coordinate Jacobian is singular,
but the coefficient-space theorem is unchanged. If a root-coordinate path
has its first nonzero imaginary coefficient at order $r$, its width
excess has first nonzero order $2r$. This explains the high even powers
seen in asymmetric frequency splitting, and extends that explanation from
one coalesced frequency to every strictly positive mirror configuration.
\end{remark}

\begin{corollary}[The quartic interaction of damping and frequency]
\label{cor:quartic}
For two targets
$s_\pm=\pm\beta+i(\kappa\pm\delta)$, with
$\abs\beta+\abs\delta$ sufficiently small relative to $\kappa$, the
shortest positive filter has width
\[
 D=\frac{2\pi}{\kappa}
 +\frac{4\pi^3}{3\kappa^5}\beta^2\delta^2
 \left[1+O\left(\frac{\beta^2+\delta^2}{\kappa^2}\right)\right].
\]
Within this neighborhood, equality $D=2\pi/\kappa$ holds exactly when
$\beta\delta=0$.
\end{corollary}

\begin{proof}
The centered polynomial is
$P(y)=y^2-(\beta+i\delta)^2$, so its normal coordinate is
$O=-2\beta\delta$. At the base $P_0(y)=y^2$ and $A(z)=(1+z)^2$.
For $V=1$, congruence \eqref{eq:hermite} gives
$B_V(z)=-h^2z/(4\kappa)$; hence the only internal velocity is
$\sigma_1=-h^2/(2\kappa)$. Formula~\eqref{eq:energy} gives
\[
 Q(1)=\frac{\sigma_1^2}{h}
       \int_0^\infty\frac{2(1+z^2)}{(1+z)^4}\dd z
      =\frac{4\sigma_1^2}{3h}=\frac{\pi^3}{3\kappa^5}.
\]
Use the analytic factorization \eqref{eq:factor}, noting that both the
real and imaginary coefficient displacements have size
$O(\beta^2+\delta^2)$ after the appropriate scaling. The exact equality
statement is Theorem~\ref{thm:main}(ii).
\end{proof}

\section{An exact example with a negative mixed curvature}

Let $\kappa=1$, $h=\pi$, and $a=\log 2/\pi$. Choose the centered targets
\[
 y_1=-2a,\quad y_2=-a,\quad y_3=3a.
\]
The actual targets $s_j=i+y_j$ all have frequency one but different
real parts. Their mean is $i$, and
\begin{align*}
 P_0(y)&=y^3-7a^2y-6a^3,\\
 A(z)&=(z+\tfrac14)(z+\tfrac12)(z+8)
      =z^3+\tfrac{35}{4}z^2+\tfrac{49}{8}z+1.
\end{align*}
All coefficients are positive. The base optimizer has weights
$(8,49,70,8)/135$ at $0,\pi,2\pi,3\pi$.

Perturb $P_0$ by $i\epsilon(v_2y+v_3)$. Formula~\eqref{eq:explicitR}
gives
\begin{align*}
 R_y(z)&=-\frac{z(196z+173)}{160a},&
 (\sigma_1,\sigma_2)&=
 \left(-\frac{173\pi}{980a},-\frac{7\pi}{50a}\right),\\
 R_1(z)&=-\frac{z(52z-49)}{160a^2},&
 (\sigma_1,\sigma_2)&=
 \left(\frac{\pi}{20a^2},-\frac{13\pi}{350a^2}\right).
\end{align*}
Writing $Q(v_2y+v_3)=(v_2,v_3)\mathsf Q(v_2,v_3)^{\mathsf T}$, exact
rational integration in \eqref{eq:energy} yields
\begin{align}
 \mathsf Q_{22}&=\frac\pi{a^2}
 \left(\frac{49}{200}-\frac{6646079\log2}{22785000}\right),\label{eq:q22}\\
 \mathsf Q_{23}&=\frac\pi{a^3}
 \left(-\frac{27}{200}+\frac{29441\log2}{155000}\right),\label{eq:q23}\\
 \mathsf Q_{33}&=\frac\pi{a^4}
 \left(\frac{21}{200}-\frac{21933\log2}{155000}\right).\label{eq:q33}
\end{align}
In particular $\mathsf Q_{23}<0$, already from $\log2<7/10$.
For example, the normalized mixed integral has the rational integrand
\[
 \frac{2912z^4-3460z^3-10919z^2+2912z-3460}
 {1000(z+8)^2(2z+1)^2(4z+1)^2}.
\]
Its partial fractions integrate to the bracket in \eqref{eq:q23}.
The full matrix is approximately
\[
 \mathsf Q=
 \begin{pmatrix}
 2.763297996804292&-0.977611543533047\\
 -0.977611543533047&9.170495331361458
 \end{pmatrix}.
\]
The Hessian is $2\mathsf Q$, not $\mathsf Q$.

Time reversal of a feasible positive measure sends target zeros $s$ to
$-\overline s$ and preserves width. In centered coordinates this replaces
$P(y)$ by $(-1)^N\overline{P(-\overline y)}$. Thus the reflected damping
example with base roots $2a,a,-3a$ has the opposite mixed entry in the same
ordered imaginary coefficient coordinates. Along a real coefficient path
between these bases through positive mirror configurations the mixed entry
therefore passes through zero. Universal positive definiteness does not
give a universal sign for individual mixed derivatives. Coordinate
orientations must be specified whenever such sign claims are made.

\section{A stiffness singularity at a three-frequency boundary}

One additional consequence illustrates what happens as the strict-positive
hypothesis degenerates. Set $\kappa=1$ and take the real frequencies
$1-d,1,1+d$, where $0<d<2/3$. Their centered polynomial in the $y$ variable
is $P_0(y)=y^3+d^2y$, and
\[
 A(z)=(z+1)(z^2+2\cos(\pi d)z+1)
      =1+bz+bz^2+z^3,\quad b=1+2\cos(\pi d)>0.
\]
For the constant normal direction $P_0+i\epsilon$, the internal velocities
are $(s,-s)$, where
$s=\pi(1-\cos(\pi d))/(bd^2)$. Hence
\begin{equation}\label{eq:tripleQ}
 Q_d=\frac{s^2}{\pi}\int_0^\infty
 \frac{b(1+z+z^3+z^4)+4b^2z^2}{(1+bz+bz^2+z^3)^2}\dd z.
\end{equation}
As $d\uparrow2/3$,
\begin{equation}\label{eq:blowup}
 \boxed{\quad Q_d\sim\frac{81\pi}{16(2/3-d)}.\quad}
\end{equation}

To prove this, divide the integrand in \eqref{eq:tripleQ} by $b$ and use
dominated convergence. For $b$ in a bounded positive interval,
$A(z)\ge1+z^3$, and the resulting polynomial numerator divided by
$(1+z^3)^2$ is integrable. The limiting integral is
\[
 \int_0^\infty\frac{1+z+z^3+z^4}{(1+z^3)^2}\dd z
 =\frac{4\pi}{3\sqrt3},
\]
by the beta integral. Since $b\sim\sqrt3\pi(2/3-d)$, substitution gives
\eqref{eq:blowup}. This is a limit of the Hessian coefficients: one first
takes the small perturbation limit at fixed $d$, then lets $d$ approach
the boundary. No uniform joint error estimate in the two parameters is
asserted. The divergence explains why a uniform coefficient neighborhood
cannot be inferred as an optimal mass approaches zero.

\section{Reproducibility, interpretation, and limits}

The accompanying \texttt{check\_mirror.py} performs exact symbolic checks
of the interpolation polynomials and the three integrals
\eqref{eq:q22}--\eqref{eq:q33}. It also solves the seven real moment
equations independently at 80 decimal digits for each of the coefficient
directions $(1,0)$, $(0,1)$, and $(1,1)$. The ratios of the computed width
excess to the proved quadratic prediction are:
\begin{center}
\begin{tabular}{@{}lrrr@{}}\toprule
Direction & $\epsilon=10^{-3}$ & $\epsilon=10^{-4}$ & $\epsilon=10^{-5}$\\\midrule
$(1,0)$ & 0.999706003 & 0.999970579 & 0.999997058\\
$(0,1)$ & 0.998232451 & 0.999822980 & 0.999982295\\
$(1,1)$ & 0.999298455 & 0.999929756 & 0.999992975\\\bottomrule
\end{tabular}
\end{center}
The largest moment residual in these nine runs is below $10^{-60}$.
The numerical calculations check the formula and conventions; the
certificate and analytic arguments prove global optimality and the
asymptotic, independently of numerical output.

In the repository's intended arithmetic application, a positive average on
logarithmic scales multiplies a term $x^s$ by $F_\mu(s)$. Theorems
\ref{thm:global} and \ref{thm:main} therefore control the least logarithmic
range needed for prescribed cancellations, including hypothetical zeros
away from a common vertical line. They neither assert that such zeros occur
nor prove that they do not. In engineering language, real parts encode
damping and imaginary parts encode oscillation; positive input shaping is
an established application of such cancellation constraints
\cite{SinghoseCrainSeering}. The new theorem is a structural duration and
stability result for the exact positive-filter problem.

The scope of the proof is finite target multisets, fixed $N$, and
neighborhoods of a base whose coefficient polynomial is strictly positive.
It supplies an exact computable Hessian but no universal neighborhood size
as $N$ varies or as a base coefficient tends to zero. It is a human-readable
proof with reproducible symbolic and numerical artifacts, not a completed
Lean formalization. Its claimed contribution is relative to the pinned
repository and the literature checked here; an exhaustive historical
priority assertion would require a broader review.

For formalization, the main dependency chain is: finite convex-hull
reduction; the strip identity for positive exponential sums; a confluent
remainder map; the real analytic implicit function theorem; persistence of
a finite certificate with simple and double zeros; uniform zero counting
by the argument principle; dominated differentiation; positive quadratic
forms; and the finite-dimensional analytic square-root map. The algebraic
Hermite interpolation and the exact rational example can be formalized
separately from the complex-analytic existence theorem.

\begin{thebibliography}{9}
\bibitem{ProveIt}
V.~Reshetnikov, ProveIt repository, \emph{Beyond the Riemann hypothesis for
Dirichlet $L$-functions: consequences, transfers and limits}, section
``Positive filters that cancel zeros'', commit
\texttt{cfab41369737c66ae520d179419509eb719367cd}.
\url{https://github.com/VladimirReshetnikov/ProveIt/blob/cfab41369737c66ae520d179419509eb719367cd/Analysis/GeneralizedRiemannHypothesis/Article/sections/fil.tex}.
In particular, labels \texttt{fil:thm:reflect}, \texttt{fil:thm:perturb},
\texttt{fil:lem:strip}, \texttt{fil:thm:meanbound},
\texttt{fil:thm:hessian}, and \texttt{fil:sec:open}.

\bibitem{EremenkoFryntov}
A.~Eremenko and A.~Fryntov, \emph{Remarks on the Obrechkoff inequality},
Proceedings of the American Mathematical Society \textbf{144} (2016),
703--707. Preprint: \url{https://arxiv.org/abs/1409.4849}.
Author-hosted paper:
\url{https://www.math.purdue.edu/~eremenko/dvi/fryntov6.pdf}.

\bibitem{SinghoseCrainSeering}
W.~Singhose, E.~Crain, and W.~Seering,
\emph{Convolved and simultaneous two-mode input shapers},
IEE Proceedings--Control Theory and Applications \textbf{144} (1997),
515--520. The bibliographic entry and background are also recorded in
the pinned repository's \texttt{Article/bib/fil.bib}.
\end{thebibliography}
\end{document}
